\begin{exercises}
  \recommended \item
    判断 \( \Re^2 \) 是否为下列每一组 \( W_1 \) 与 \( W_2 \) 的直和。
    \begin{exparts}
       \partsitem \( W_1=\set{\colvec{x \\ 0}\suchthat x\in\Re} \),
             \( W_2=\set{\colvec{x \\ x}\suchthat x\in\Re} \)
       \partsitem \( W_1=\set{\colvec{s \\ s}\suchthat s\in\Re} \),
             \( W_2=\set{\colvec{s \\ 1.1s}\suchthat s\in\Re} \)
       \partsitem \( W_1=\Re^2 \), \( W_2=\set{\zero} \)
       \partsitem \( W_1=W_2=\set{\colvec{t \\ t}\suchthat t\in\Re} \)
       \partsitem \( W_1=\set{\colvec[r]{1 \\ 0}+\colvec{x \\ 0}
                                \suchthat x\in\Re} \),
             \( W_2=\set{\colvec[r]{-1 \\ 0}+\colvec[r]{0 \\ y}\suchthat y\in\Re} \)
    \end{exparts}
    \begin{answer}
       对这里的每一项我们都可以应用\nearbylemma{le:DirectSumTwoSp}。
       \begin{exparts}
         \partsitem 是的。
           该平面是这个 $W_1$ 与 $W_2$ 的和，因为对任意标量 $a$ 与 $b$，
           \begin{equation*}
             \colvec{a \\ b}
             =\colvec{a-b \\ 0}
             +\colvec{b \\ b}
           \end{equation*}
           表明任一向量都是来自这两部分的向量之和。
           而且，这两个子空间是（不同的）过原点的直线，
           因此它们的交是平凡的。
         \partsitem 是的。
           为看出平面中任一向量都是这些部分中向量的组合，
           考察下面这个关系式。
           \begin{equation*}
             \colvec{a \\ b}=c_1\colvec[r]{1 \\ 1}+c_2\colvec[r]{1 \\ 1.1}
           \end{equation*}
           我们现在只需注意，集合
           \begin{equation*}
             \set{\colvec[r]{1 \\ 1},\colvec[r]{1 \\ 1.1}}
           \end{equation*}
           是该空间的一组基（因为它显然线性无关，且在 $\Re^2$ 中含有两个元素），
           因此上述方程有且仅有一解，这意味着所有分解都是唯一的。
           另一种做法是求解
           \begin{equation*}
             \begin{linsys}{2}
               c_1  &+  &c_2    &=  &a  \\
               c_1  &+  &1.1c_2 &=  &b
             \end{linsys}
             \grstep{-\rho_1+\rho_2}
             \begin{linsys}{2}
               c_1  &+  &c_2    &=  &a  \\
                    &   &0.1c_2 &=  &-a+b
             \end{linsys}
           \end{equation*}
           得到 $c_2=10(-a+b)$ 与 $c_1=11a-10b$，于是我们有
           \begin{equation*}
             \colvec{a \\ b}=\colvec{11a-10b \\ 11a-10b}
                             +\colvec{-10a+10b \\ 1.1\cdot(-10a+10b)}
           \end{equation*}
           正如所求。
           与前一个答案一样，这两个子空间各自都是过原点的直线，
           它们的交是平凡的。
         \partsitem 是的。
           平面中的每个向量都可如下分解为和
           \begin{equation*}
             \colvec{x \\ y}=\colvec{x \\ y}+\colvec[r]{0 \\ 0}
           \end{equation*}
           且这两个子空间的交是平凡的。
         \partsitem 不是。
           交不是平凡的。
         \partsitem 不是。
           它们不是子空间。
      \end{exparts}  
    \end{answer}
  \recommended \item
    证明 \( \Re^3 \) 是 \( xy \) 平面与下列每一个的直和。
    \begin{exparts}
      \partsitem \( z \) 轴
      \partsitem 直线
        \begin{equation*}
          \set{\colvec{z \\ z \\ z} \suchthat z\in\Re }
        \end{equation*}
    \end{exparts}
    \begin{answer}
      对这里的每一项我们都可以使用\nearbylemma{le:DirectSumTwoSp}。      
      \begin{exparts}
        \partsitem \( \Re^3 \) 中的任意向量都可以分解为这个和。
          \begin{equation*}
            \colvec{x \\ y \\ z}
            =\colvec{x \\ y \\ 0}
            +\colvec{0 \\ 0 \\ z}
          \end{equation*}
          而且，$xy$ 平面与 $z$ 轴的交是平凡子空间。
        \partsitem \( \Re^3 \) 中的任意向量都可以分解为
          \begin{equation*}
            \colvec{x \\ y \\ z}
            =\colvec{x-z \\ y-z \\ 0}
            +\colvec{z \\ z \\ z}
          \end{equation*}
          且这两个空间的交是平凡的。
      \end{exparts}  
    \end{answer}
  \item 
    \( \polyspace_2 \) 是
    \( \set{a+bx^2\suchthat a,b\in\Re} \) 与
    \( \set{cx\suchthat c\in\Re} \) 的直和吗？
    \begin{answer}
      是。
      证明这两个集合是子空间是常规的。
      为看出整个空间是这两个的直和，只需注意
      \( \polyspace_2 \) 中每个元素都有唯一的分解
      \( m+nx+px^2=(m+px^2)+(nx) \)。  
     \end{answer}
  \recommended \item 
    在 \( \polyspace_n \) 中，
    \definend{偶}\index{even functions}多项式是这个集合中的元素
    \begin{equation*}
      \mathcal{E}=
      \set{p\in\polyspace_n \suchthat \text{\( p(-x)=p(x) \) 对一切 \( x \)}}
    \end{equation*}
    而
    \definend{奇}\index{odd function}
    多项式是这个集合中的元素。
    \begin{equation*}
      \mathcal{O}=
      \set{p\in\polyspace_n \suchthat \text{\( p(-x)=-p(x) \) 
                                              对一切 \( x \)}}
    \end{equation*}
    证明它们是互补的子空间。
    \begin{answer}
      证明它们是子空间是常规的。
      我们将利用\nearbylemma{le:DirectSumTwoSp}来论证它们互补。
      交 \( \mathcal{E}\intersection\mathcal{O} \) 是平凡的，
      因为同时满足条件 $p(-x)=p(x)$ 与
      $p(-x)=-p(x)$ 的多项式只有零多项式。
      为看出整个空间是这些子空间的和
      \( \mathcal{E}+\mathcal{O}=\polyspace_n \)，
      注意多项式
      \( p_0(x)=1 \)、\( p_2(x)=x^2 \)、\( p_4(x)=x^4 \) 等都在
      \( \mathcal{E} \) 中，同时注意多项式
      \( p_1(x)=x \)、\( p_3(x)=x^3 \) 等都在 \( \mathcal{O} \) 中。
      于是 \( \polyspace_n \) 的任何元素都是
      \( \mathcal{E} \) 与 \( \mathcal{O} \) 中元素的组合。  
    \end{answer}
  \item 
    下列 \( \Re^3 \) 的哪些子空间
    \begin{center}
     \begin{tabular}{l}
      $W_1$: \( x \) 轴，\quad
      $W_2$:~\( y \) 轴，\quad
      $W_3$:~\( z \) 轴，             \\
      $W_4$:~平面 \( x+y+z=0 \)，\quad
      $W_5$:~\( yz \) 平面
     \end{tabular}
    \end{center}
    可以组合成
    \begin{exparts*}
      \partsitem 和为 \( \Re^3 \)？
      \partsitem 直和为 \( \Re^3 \)？
    \end{exparts*}
    \begin{answer}
     这里的每一个都是 $\Re^3$。
     \begin{exparts}
      \partsitem 
        为便于阅读，这里把它们分成若干行列出。
        \begin{center}
          \begin{tabular}{l} 
            $W_1+W_2+W_3$, 
              $W_1+W_2+W_3+W_4$, 
              $W_1+W_2+W_3+W_5$,  \\ \quad %line too long
              $W_1+W_2+W_3+W_4+W_5$, 
              $W_1+W_2+W_4$, 
              $W_1+W_2+W_4+W_5$,     \\ \quad %line too long
              $W_1+W_2+W_5$,         
              $W_1+W_3+W_4$, 
              $W_1+W_3+W_5$, 
              $W_1+W_3+W_4+W_5$,     \\ \quad 
              $W_1+W_4$, 
              $W_1+W_4+W_5$,         
              $W_1+W_5$,             \\ 
              $W_2+W_3+W_4$, 
              $W_2+W_3+W_4+W_5$,     
              $W_2+W_4$, 
              $W_2+W_4+W_5$,         \\ 
              $W_3+W_4$, 
              $W_3+W_4+W_5$,         \\ 
              $W_4+W_5$ 
          \end{tabular}
         \end{center}
     \partsitem
       $W_1\directsum W_2\directsum W_3$,    
       $W_1\directsum W_4$,    
       $W_1\directsum W_5$,    
       $W_2\directsum W_4$,    
       $W_3\directsum W_4$    
     \end{exparts}
    \end{answer}
  \recommended \item
    证明 \( \polyspace_n=\set{a_0 \suchthat a_0\in\Re}\directsum\dots
                              \directsum\set{a_nx^n\suchthat a_n\in\Re} \)。
    \begin{answer}
      显然每一个都是子空间。
      这些子空间的基 \( B_i=\sequence{x^i} \) 依次连接起来，
      就构成整个空间的一组基。  
    \end{answer}
  \item 
    若 \( W_1\subseteq W_2 \)，则 \( W_1+W_2 \) 是什么？
    \begin{answer}
       是 \( W_2 \)。  
    \end{answer}
  \item 
    \nearbyexample{exam:BenchDirSum} 能否推广？
    也就是说，下述论断是真是假：~若向量空间 \( V \) 有一组基
    \( \sequence{\vec{\beta}_1,\dots ,\vec{\beta}_n} \)，则
    它是这些一维子空间之张成的直和
    \( V=\spanof{\set{\vec{\beta}_1}}\directsum\dots
           \directsum\spanof{\set{\vec{\beta}_n}} \)？
    \begin{answer}
      由\nearbylemma{le:UniqDecIffBasisDec}知其为真。  
    \end{answer}
  \item 
    \( \Re^4 \) 能以两种不同的方式分解为直和吗？
    \( \Re^1 \) 呢？
    \begin{answer}
      \( \Re^4 \) 的两个不同的直和分解很容易找到。
      其一是
      \( W_1=\spanof{\set{\vec{e}_1,\vec{e}_2}} \) 与
      \( W_2=\spanof{\set{\vec{e}_3,\vec{e}_4}} \)，
      其二是
      \( U_1=\spanof{\set{\vec{e}_1}} \) 与
      \( U_2=\spanof{\set{\vec{e}_2,\vec{e}_3,\vec{e}_4}} \)。
      （还有很多种可能，
      例如 \( \Re^4 \) 与它的平凡子空间。）

      相比之下，对 \( \Re^1 \) 的单向量基的任何划分，
      都会给出一组没有元素的基和另一组只有一个元素的基。
      因此任何分解都涉及 \( \Re^1 \) 与它的平凡
      子空间。 
    \end{answer}
  \item 
      本练习使在集合之间写 `$+$' 的记法显得更为自然。 
      证明：设 \( W_1,\dots, W_k \) 是某个向量空间的子空间，则
      \begin{equation*}
        W_1+\dots+W_k
        =\set{\vec{w}_1+\vec{w}_2+\dots+\vec{w}_k
              \suchthat \vec{w}_1\in W_1,\dots,\vec{w}_k\in W_k},
      \end{equation*}
      从而子空间的和就是由所有和构成的子空间。
      \begin{answer}
        一个方向的包含关系是容易的：
        \( \set{\vec{w}_1+\dots+\vec{w}_k\suchthat \vec{w}_i\in W_i} \)
        是 \( \spanof{W_1\union \dots \union W_k}  \) 的子集，
        因为每个 \( \vec{w}_1+\dots+\vec{w}_k \) 都是这个并集中向量的和。

        对于另一个方向的包含，对来自该并集的向量的任一线性组合，
        利用向量加法的交换性，
        把来自 \( W_1 \) 的向量放在最前面，
        接着是来自 \( W_2 \) 的向量，等等。
        把来自 \( W_1 \) 的向量相加得到一个 \( \vec{w}_1\in W_1 \)，
        把来自 \( W_2 \) 的向量相加得到一个 \( \vec{w}_2\in W_2 \)，等等。
        所得结果具有所要求的形式。
      \end{answer}
  \item \label{exer:ThreeSubsPairwseNonTriv}
    （参见\nearbyexample{exam:DirSumThree}。
    本练习
    表明，“两两交都平凡”这一要求
    确实强于只要求所有子空间的交为平凡。）
    给出一个向量空间和三个子空间 $W_1$、$W_2$ 与 $W_3$，
    使得该空间是这些子空间的和，
    三个子空间的交
    $W_1\intersection W_2\intersection W_3$ 是平凡的，
    但两两的交
    $W_1\intersection W_2$、$W_1\intersection W_3$ 与 $W_2\intersection W_3$
    都不是平凡的。
    \begin{answer}
      一个例子是取空间为 $\Re^3$，
      并取子空间为 $xy$ 平面、$xz$ 平面与 $yz$ 平面。
    \end{answer}
  \item 
    证明：若 \( V=W_1\directsum\dots\directsum W_k \)，则
    只要 \( i\neq j \)，\( W_i\intersection W_j \) 就是平凡的。
    这表明\nearbylemma{le:DirectSumTwoSp}
    证明的前一半可以推广到多于两个
    子空间的情形。
    （\nearbyexample{exam:DirSumThree} 表明这个蕴含关系
    不能反向；后一半不能推广。）
    \begin{answer}
      当然，零向量在所有子空间中，所以
      交至少包含这一个向量。  
      由直和的定义，集合
      \( \set{W_1,\dots,W_k} \) 是无关的，因此当
      \( i\neq j \) 时，\( W_i \) 中没有非零向量
      是 \( W_j \) 中某个元素的倍数。
      特别是，\( W_i \) 中没有非零向量等于
      \( W_j \) 中的元素。
    \end{answer}
  \item 
    回想一下，任何线性无关集合都不包含零向量。
    一个由子空间组成的无关集合能包含平凡子空间吗？
    \begin{answer}
      它可以包含一个平凡子空间；
      \( \Re^3 \) 的下面这个子空间集合是无关的：
      \( \set{\set{\zero},\text{\( x \)-轴}} \)。
      平凡空间 \( \set{\zero} \) 中没有非零向量
      是 \( x \) 轴中某个
      向量的倍数，这仅仅是因为平凡空间
      没有可以作为这种倍数候选者的非零向量（而且
      \( x \) 轴中也没有非零向量是平凡子空间中
      零向量的倍数）。  
     \end{answer}
  \recommended \item 
    每个子空间都有补吗？
    \begin{answer}
      有。
      对向量空间的任何子空间，我们可以取该
      子空间的任意一组基 \( \sequence{\vec{\omega}_1,\dots,\vec{\omega}_k} \)，
      并把它扩充成整个
      空间的一组基
      \( \sequence{\vec{\omega}_1,\dots,\vec{\omega}_k,
                     \vec{\beta}_{k+1},\dots,\vec{\beta}_n} \)。
      于是原子空间的补具有这组基
      \( \sequence{\vec{\beta}_{k+1},\dots,\vec{\beta}_n} \)。  
     \end{answer}
  \recommended \item 
    设 \( W_1, W_2 \) 是某个向量空间的子空间。
    \begin{exparts}
      \partsitem 假设
        集合 \( S_1 \) 张成 \( W_1 \)，且集合 \( S_2 \) 张成
        \( W_2 \)。
        \( S_1\union S_2 \) 能张成 \( W_1+W_2 \) 吗？
        它必定张成吗？
      \partsitem 假设 \( S_1 \) 是 \( W_1 \) 的一个线性无关
        子集，
        且 \( S_2 \) 是 \( W_2 \) 的一个线性无关子集。
        \( S_1\union S_2 \) 能是
        \( W_1+W_2 \) 的线性无关子集吗？
        它必定是吗？
    \end{exparts}
    \begin{answer}
      \begin{exparts}
        \partsitem 它必定张成。
          我们可以把 \( W_1+W_2 \) 的任何成员写成 \( \vec{w}_1+\vec{w}_2 \)，
          其中 \( \vec{w}_1\in W_1 \) 且 \( \vec{w}_2\in W_2 \)。
          由于 \( S_1 \) 张成 \( W_1 \)，向量 \( \vec{w}_1 \) 是
          \( S_1 \) 中成员的组合。
          类似地，\( \vec{w}_2 \) 是 \( S_2 \) 中成员的组合。
        \partsitem 看出它可以是线性无关的一个简单方法
          是取两者都为空集。
          另一方面，在空间 $\Re^1$ 中，
          若 \( W_1=\Re^1 \) 且 \( W_2=\Re^1 \)，而 $S_1=\set{1}$，
          $S_2=\set{2}$，则它们的并 $S_1\union S_2$ 不是无关的。
      \end{exparts}  
     \end{answer}
  \item \label{exer:BasesSumTwoSubs}
    当我们把一个向量空间分解为直和时，各子空间的
    维数相加等于空间的维数。
    而当一个空间是作为其子空间的和给出时，情形
    就没有这么简单。
    本练习考虑两个子空间的特殊情形。
    \begin{exparts}
      \partsitem 对 \( \matspace_{\nbyn{2}} \) 的这些子空间，求
        \( W_1\intersection W_2 \)、\( \dim(W_1\intersection W_2) \)、
        \( W_1+W_2 \) 与 \( \dim(W_1+W_2) \)。
        \begin{equation*}
          W_1=\set{\begin{mat}
                    0  &0  \\
                    c  &d
                  \end{mat} \suchthat c,d\in\Re  }
          \qquad
         W_2=\set{\begin{mat}
                    0  &b  \\
                    c  &0
                  \end{mat} \suchthat b,c\in\Re  }
       \end{equation*}
     \partsitem 设 \( U \) 与 \( W \) 是某个向量空间的
       子空间。
       设序列
       \( \sequence{\vec{\beta}_1,\dots,\vec{\beta}_k} \)
       是
       \( U\intersection W \) 的一组基。
       最后，设前述序列已被扩充，得到
       \( U \) 的一组基 \( \sequence{\vec{\mu}_1,\dots,\vec{\mu}_j,
       \vec{\beta}_1,\dots,\vec{\beta}_k} \)，
       以及 \( W \) 的一组基
       \( \sequence{\vec{\beta}_1,\dots,\vec{\beta}_k,
          \vec{\omega}_1,\dots,\vec{\omega}_p} \)。
       证明序列
       \begin{equation*}
         \sequence{\vec{\mu}_1,\dots,\vec{\mu}_j,
              \vec{\beta}_1,\dots,\vec{\beta}_k,
              \vec{\omega}_1,\dots,\vec{\omega}_p}
       \end{equation*}
       是和 \( U+W \) 的一组基。
      \partsitem 由此得出结论
          $\dim (U+W)=\dim(U)+\dim(W)-\dim(U\intersection W)$。
      \partsitem 设 \( W_1 \) 与 \( W_2 \) 是十维空间的两个八维
        子空间。
       列出 \( \dim(W_1\intersection W_2) \) 的所有可能取值。
    \end{exparts}
    \begin{answer}
      \begin{exparts}
        \partsitem 交与和是
          \begin{equation*}
             \set{\begin{mat}
                    0  &0  \\
                    c  &0
                  \end{mat} \suchthat c\in\Re  }
             \qquad
             \set{\begin{mat}
                    0  &b  \\
                    c  &d
                   \end{mat} \suchthat b,c,d\in\Re  }
          \end{equation*}
         其维数分别为一与三。
      \partsitem 我们把 $U\intersection W$ 的基记为
        $B_{U\intersection W}$，
        把 $U$ 的基记为 $B_U$，
        把 $W$ 的基记为 $B_W$，
        而把正在考虑的基记为 $B_{U+W}$。

        为看出 $B_{U+W}$ 张成 $U+W$，注意我们可以把
        $U+W$ 中的任一向量 $c\vec{u}+d\vec{w}$ 写成
        $B_{U+W}$ 中向量的线性
        组合，
        只需把 $\vec{u}$ 用
        $B_U$ 表示、把 $\vec{w}$ 用 $B_W$ 表示即可。

         我们最后证明 $B_{U+W}$ 线性无关。 
         考察
         \begin{equation*}
           c_1\vec{\mu}_1+\dots+c_{j+1}\vec{\beta}_1+\dots
              +c_{j+k+p}\vec{\omega}_p=\zero
         \end{equation*}
         它可以如下改写。
         \begin{equation*}
           c_1\vec{\mu}_1+\dots+c_j\vec{\mu}_j=
             -c_{j+1}\vec{\beta}_1-\dots-c_{j+k+p}\vec{\omega}_p
         \end{equation*}
         注意左边求和得到 \( U \) 中的向量，而
         右边求和得到 \( W \) 中的向量，于是两边求和都得到
         $U\intersection W$ 的一个成员。
         由于左边是 $U\intersection W$ 的成员，
         它可以用 $B_{U\intersection W}$ 的成员表示，
         这给出：上式左边 $\vec{\mu}$ 的组合
         等于 $\vec{\beta}$ 的一个组合。 
         但是，基 $B_U$ 线性无关这一事实表明
         任何这样的组合
         都是平凡的，特别地，上面左边的系数 $c_1$, \ldots,
         $c_j$ 全为零。
         类似地，\( \vec{\omega} \) 的系数也全为零。
         这使上面的方程成为
         \( \vec{\beta} \) 之间的一个线性关系，
         但 $B_{U\intersection W}$ 线性无关，
         因此 $\vec{\beta}$ 的所有系数也都
         为零。
        \partsitem 只需数一数上一项中的
          基向量：~$\dim(U+W)=j+k+p$，而 $\dim(U)=j+k$，$\dim(W)=k+p$，
          以及 $\dim(U\intersection W)=k$。
        \partsitem  我们知道
          \( \dim(W_1+W_2)=\dim(W_1)+\dim(W_2)-\dim(W_1\intersection W_2) \)。
          因为 \( W_1\subseteq W_1+W_2 \)，
          我们知道 \( W_1+W_2 \) 的维数必定大于
          $W_1$ 的维数，也就是说，其维数必定是八、九或十。
          代入得到三种可能：
          $8=8+8-\dim(W_1\intersection W_2)$、
          $9=8+8-\dim(W_1\intersection W_2)$ 或
          $10=8+8-\dim(W_1\intersection W_2)$。
          因此 \( \dim(W_1\intersection W_2) \) 必定是
          八、七或六。
          （要举例说明这三种情形中的每一种都可能
          出现是容易的，例如在 \( \Re^{10} \) 中。）  
      \end{exparts}
    \end{answer}
  \item 
    设 \( V=W_1\directsum\cdots\directsum W_k \)，并对每个下标
    \( i \)，设 \( S_i \) 是
    \( W_i \) 的一个线性无关子集。
    证明诸 \( S_i \) 的并集线性无关。
    \begin{answer}
      把每个 \( S_i \) 扩充成 \( W_i \) 的一组基 $B_i$。
      这些基的连接 $\cat{\cat{B_1}{\cdots}}{B_k}$
      是 \( V \) 的一组基，因此其中的成员构成一个线性无关
      集合。
      但并集 $S_1\union\cdots\union S_k$ 是这个线性无关集合的
      子集，因此它自身
      也线性无关。  
     \end{answer}
  \item 
    一个矩阵是\definend{对称的}\index{matrix!symmetric}%
    \index{symmetric matrix}，若对每一对下标 \( i \) 与
    \( j \)，\( i,j \) 元等于 \( j,i \) 元。
    一个矩阵是\definend{反对称的}\index{matrix!antisymmetric}%
    \index{antisymmetric matrix}，若每个 \( i,j \) 元是
    \( j,i \) 元的负值。
    \begin{exparts}
      \partsitem 给出一个对称的 $\nbyn{2}$ 矩阵和一个反对称的
        $\nbyn{2}$ 矩阵。
        （\textit{注}：
        对于第二个，要当心对角线上的元素。）
      \partsitem 对称方阵与它的转置之间有什么关系？
        反对称方阵与它的转置之间呢？
      \partsitem 证明 \( \matspace_{\nbyn{n}} \) 是对称矩阵空间与反对称
        矩阵空间的直和。
    \end{exparts}
    \begin{answer}
      \begin{exparts}
        \partsitem 这样的两个矩阵如下。
          \begin{equation*}
            \begin{mat}[r]
              1  &2  \\
              2  &3
            \end{mat}
            \qquad
            \begin{mat}[r]
              0  &1  \\
              -1 &0
            \end{mat}
          \end{equation*}
          对于反对称的那个，对角线上的元素必须为零。
        \partsitem 对称方阵等于它的转置。
          反对称方阵等于它的转置的负值。
        \partsitem 证明这两个集合是子空间是容易的。
         设 \( A\in\matspace_{\nbyn{n}} \)。
         为把 $A$ 表示成一个对称矩阵与一个反对称矩阵的和，
         我们注意到
         \begin{equation*}
            A=(1/2)(A+\trans{A}) + (1/2)(A-\trans{A})
         \end{equation*}
         并注意第一个加项是对称的，而第二个是
         反对称的。
         因此 \( \matspace_{\nbyn{n}} \) 是这两个子空间的和。
         为证明这个和是直和，
         设矩阵 \( A \) 既是对称的
         \( A=\trans{A} \) 又是反对称的 \( A=-\trans{A} \)。
         则 \( A=-A \)，于是 \( A \) 的所有元素都是零。  
      \end{exparts}
     \end{answer}
  \item 
    设 \( W_1,W_2,W_3 \) 是某个向量空间的子空间。
    证明 \( (W_1\intersection W_2)+(W_1\intersection W_3)\subseteq
              W_1\intersection (W_2+W_3) \)。
    这个包含关系能反向吗？
    \begin{answer}
      设
      \( \vec{v}\in (W_1\intersection W_2)+(W_1\intersection W_3) \)。
      则 \( \vec{v}=\vec{w}_2+\vec{w}_3 \)，其中
      \( \vec{w}_2\in W_1\intersection W_2 \) 且
      \( \vec{w}_3\in W_1\intersection W_3 \)。
      注意 \( \vec{w}_2,\vec{w}_3\in W_1 \)，又子空间对加法
      封闭，故 \( \vec{w}_2+\vec{w}_3\in W_1 \)。
      于是 \( \vec{v}=\vec{w}_2+\vec{w}_3\in W_1\intersection(W_2+W_3) \)。

      这个例子证明包含可能是严格的：
      在 \( \Re^2 \) 中取 \( W_1 \) 为 \( x \) 轴，取 \( W_2 \)
      为 \( y \) 轴，取 \( W_3 \) 为直线 \( y=x \)。
      则 \( W_1\intersection W_2 \) 与 \( W_1\intersection W_3 \)
      都是平凡的，因此它们的和是平凡的。
      但 \( W_2+W_3 \) 是整个 \( \Re^2 \)，所以
      \( W_1\intersection(W_2+W_3) \) 是 \( x \) 轴。  
     \end{answer}
  \item 
    \( \Re^2 \) 中 $x$ 轴与 $y$ 轴的例子
    表明 \( W_1\directsum W_2=V \) 并不蕴含
    \( W_1\union W_2=V \)。
    \( W_1\directsum W_2=V \) 与 \( W_1\union W_2=V \)
    能同时发生吗？
    \begin{answer}
      当 \( W_1,W_2 \) 中至少有一个是平凡的时，这会发生。
      但这是它可能发生的唯一方式。

      为证明这一点，设两者都是非平凡的，
      从各自中取非零向量 \( \vec{w}_1, \vec{w}_2 \)，
      并考察 \( \vec{w}_1+\vec{w}_2 \)。
      这个和不属于 \( W_1 \)，因为
      \( \vec{w}_1+\vec{w}_2=\vec{v}\in W_1 \) 将蕴含
      \( \vec{w}_2=\vec{v}-\vec{w}_1 \) 属于 \( W_1 \)，
      这与子空间无关性的假设矛盾。
      类似地，\( \vec{w}_1+\vec{w}_2 \) 不属于 \( W_2 \)。
      于是 \( V \) 中有一个元素不属于 \( W_1\union W_2 \)。
     \end{answer}
%   \recommended \item
%     Our model for complementary subspaces, the $x$-axis and the
%     $y$-axis in \( \Re^2 \), has one property not used here.
%     Where \( U \) is a subspace of \( \Re^n \) we define the
%     \definend{orthocomplement}\index{subspace!orthocomplement} 
%     of \( U \) to be
%     \begin{equation*}
%       U^\perp=
%       \set{\vec{v}\in\Re^n \suchthat \text{\( \vec{v}\dotprod\vec{u}=0 \)
%                                             for all \( \vec{u}\in U \)}}
%     \end{equation*}
%     (read ``\( U \) perp'').
%     \begin{exparts}
%       \partsitem Find the orthocomplement of the \( x \)-axis in \( \Re^2 \).
%       \partsitem Find the orthocomplement of the \( x \)-axis in \( \Re^3 \).
%       \partsitem Find the orthocomplement of the \( xy \)-plane in \( \Re^3 \).
%       \partsitem Show that the orthocomplement of a subspace is a subspace.
%       \partsitem Show that if \( W \) is the orthocomplement of \( U \) then
%         \( U \) is the orthocomplement of \( W \).
%       \partsitem Prove that a subspace and its orthocomplement have a trivial
%         intersection.
%       \partsitem Conclude that for any \( n \) and subspace 
%         \( U\subseteq \Re^n \) we have that \( \Re^n=U\directsum U^\perp \).
%       \partsitem Show that 
%         \( \dim(U)+\dim(U^\perp) \) equals the dimension of the
%         enclosing space.
%     \end{exparts}
%     \begin{answer}
%       \begin{exparts}
%         \partsitem The set 
%           \begin{equation*}
%             \set{\colvec{v_1 \\ v_2} 
%                \suchthat \colvec{v_1 \\ v_2}\dotprod\colvec{x \\ 0}=0 
%                               \text{\ for all $x\in\Re$}}
%           \end{equation*}
%           is easily seen to be the $y$-axis.
%         \partsitem The \( yz \)-plane.
%         \partsitem The \( z \)-axis.
%         \partsitem 
%           Assume that \( U \) is a subspace of some \( \Re^n \).
%           Because $U^\perp$ contains the zero vector, since that
%           vector is perpendicular to everything, we need only show that the
%           orthocomplement is closed under linear combinations of two elements.
%           If \( \vec{w}_1, \vec{w}_2\in U^\perp \) then
%           \( \vec{w}_1\dotprod\vec{u}=0 \) and \( \vec{w}_2\dotprod\vec{u}=0 \)
%           for all \( \vec{u}\in U \).
%           Thus \( (c_1\vec{w}_1+c_2\vec{w}_2)\dotprod\vec{u}=
%               c_1(\vec{w}_1\dotprod\vec{u})+c_2(\vec{w}_2\dotprod\vec{u})=0 \)
%           for all \( \vec{u}\in U \) and so \( U^\perp \) is closed under
%           linear combinations.
%         \partsitem The only vector orthogonal to itself is the zero vector.
%         \partsitem This is immediate.
%         \partsitem To prove that the dimensions add, 
%           it suffices by \nearbycorollary{cor:DirSumDimsAdd} and 
%           \nearbylemma{le:DirectSumTwoSp} to show that 
%           \( U\intersection U^\perp \) is the trivial subspace
%           \( \set{\zero} \).
%           But this is one of the prior items in this problem.
%       \end{exparts}  
%      \end{answer}
  \item 
    考察\nearbycorollary{cor:DirSumDimsAdd}。
    它反过来也成立吗\Dash 也就是说，
    设 \( V=W_1+\dots+ W_k \)，
    是否 \( V=W_1\directsum\cdots\directsum W_k \) 当且仅当
    \( \dim(V)=\dim(W_1)+\dots+\dim(W_k) \)？
    \begin{answer}
      是。
      从左到右的蕴含就是
      \nearbycorollary{cor:DirSumDimsAdd}。
      对于另一个方向，
      设 \( \dim(V)=\dim(W_1)+\dots+\dim(W_k) \)。
      取 \( B_1,\dots, B_k \) 为 \( W_1,\dots, W_k \) 的基。
      由于 $V$ 是这些子空间的和，我们可以把
      任一 \( \vec{v}\in V \) 写成
      \( \vec{v}=\vec{w}_1+\cdots+\vec{w}_k \)，而把每个
      \( \vec{w}_i \) 表示成相应
      基 \( B_i \) 中向量的组合，即表明连接
      \( \cat{B_1}{\cat{\cdots}{B_k}}  \) 张成 \( V \)。
      现在，这个连接有
      \( \dim(W_1)+\dots+\dim(W_k) \) 个成员，因此它是一个大小为
      \( \dim(V) \) 的张成集。
      所以这个连接是 \( V \) 的一组基。
      于是 \( V \) 是直和。
     \end{answer}
  \item 
    我们知道，若 \( V=W_1\directsum W_2 \)，则 \( V \) 有一组基
    可以拆分成 \( W_1 \) 的一组基与
    \( W_2 \) 的一组基。
    我们能作出更强的论断，说 \( V \) 的每一组基都能拆分成
    \( W_1 \) 的一组基与 \( W_2 \) 的一组基吗？
    \begin{answer}
      不能。
      \( \Re^2 \) 的标准基不能拆分成互补子空间，即直线 \( x=y \) 与直线
      \( x=-y \) 的基。  
    \end{answer}
  \item 
     我们可以问一问 `$+$' 运算的代数性质。
     \begin{exparts}
       \partsitem 它是否交换；即 \( W_1+W_2=W_2+W_1 \) 是否成立？
       \partsitem 它是否结合；即 \( (W_1+W_2)+W_3=W_1+(W_2+W_3) \)
         是否成立？
       \partsitem 设 \( W \) 是某个向量空间的子空间。
         证明 \( W+W=W \)。
       \partsitem 必须存在单位元吗？
         也就是一个子空间 \( I \)，使得 \( I+W=W+I=W \)
         对一切子空间 \( W \) 成立？ 
       \partsitem 左消去成立吗：~若
          \( W_1+W_2=W_1+W_3 \)，则 \( W_2=W_3 \)？
          右消去呢？
     \end{exparts}
    \begin{answer}
      \begin{exparts}
        \partsitem  是的，对一切子空间 \( W_1,W_2 \)
          都有 \( W_1+W_2=W_2+W_1 \)，
          因为每一边都是 \( W_1\union W_2=W_2\union W_1 \) 的张成。
        \partsitem 这一条与前一条类似\Dash 该等式的每一边
          都是
          \( (W_1\union W_2)\union W_3=W_1\union (W_2\union W_3) \)
          的张成。
        \partsitem 由于这是集合之间的等式，我们可以用相互包含
          来证明它成立。
          显然 \( W\subseteq W+W \)。
          对于 \( W+W\subseteq W \)，只需记得每个子集都对加法封闭，
          所以形如 \( \vec{w}_1+\vec{w}_2 \) 的任何和都在
          \( W \) 中。
        \partsitem 在每个向量空间中，关于子空间加法的
          单位元是平凡子空间。
        \partsitem   左消去与右消去都不必成立。
          举个例子，在 \( \Re^3 \) 中取 \( W_1 \) 为
          \( xy \) 平面，取 \( W_2 \) 为 \( x \) 轴，
          取 \( W_3 \) 为 \( y \) 轴。  
      \end{exparts}
     \end{answer}
  % 2014-Dec-19 I commented this out to make the pages fit better; it is perfectly fine
  % \item 
  %   Consider the algebraic properties of the direct sum operation.
  %   \begin{exparts}
  %      \partsitem Does direct sum commute: does \( V=W_1\directsum W_2 \) imply
  %        that \( V=W_2\directsum W_1 \)?
  %     \partsitem Prove that direct sum is associative:
  %       \( (W_1\directsum W_2)\directsum W_3=
  %           W_1\directsum(W_2\directsum W_3) \).
  %      \partsitem Show that \( \Re^3 \) is the direct sum of the three axes
  %        (the relevance here is that by the previous item,
  %         we needn't specify which two of the three axes are combined first).
  %      \partsitem Does the direct sum operation left-cancel:~does
  %        \( W_1\directsum W_2=W_1\directsum W_3 \) imply \( W_2=W_3 \)?
  %        Does it right-cancel?
  %      \partsitem There is an identity element with respect to this operation.
  %        Find it.
  %      \partsitem Do some, or all, subspaces have inverses with respect to this
  %        operation:~is there a subspace \( W \) of some vector space such
  %        that there is a subspace \( U \) with the property that
  %        \( U\directsum W \) equals the identity element from
  %        the prior item?
  %   \end{exparts}
  %   \begin{answer}
  %     \begin{exparts}
  %        \partsitem They are equal because for each, 
  %          \( V \) is the direct sum if
  %          and only if we can write each \( \vec{v}\in V \) in a unique
  %          way as a sum \( \vec{v}=\vec{w}_1+\vec{w}_2 \) and
  %          \( \vec{v}=\vec{w}_2+\vec{w}_1 \).
  %        \partsitem They are equal because for each, 
  %          \( V \) is the direct sum if
  %          and only if we can write each \( \vec{v}\in V \) in a unique
  %          way as a sum of a vector from each
  %          $\vec{v}=(\vec{w}_1+\vec{w}_2)+\vec{w}_3$
  %          and $\vec{v}=\vec{w}_1+(\vec{w}_2+\vec{w}_3)$.
  %        \partsitem We can decompose any vector in \( \Re^3 \) uniquely into
  %          the sum of a vector from each axis.
  %        \partsitem No.
  %          For an example, in \( \Re^2 \) take \( W_1 \) to be the
  %          \( x \)-axis, take \( W_2 \) to be the \( y \)-axis, and
  %          take \( W_3 \) to be the line \( y=x \).
  %        \partsitem In any vector space the trivial subspace acts as 
  %          the identity element with respect to direct sum.
  %        \partsitem In any vector space, only the trivial subspace has
  %          a direct-sum inverse (namely, itself).
  %          One way to see this is that dimensions add, and so increase.
  %     \end{exparts}  
  %    \end{answer}
\end{exercises}
\index{direct sum|)}
